\documentclass[sigconf,nonacm,review,natbib=true]{acmart}

%% Inconsolata has no italic face in the bundled AAMAS toolchain.  Listings
%% occasionally request one through a language style, so make the intended
%% upright substitution explicit rather than leaving a font fallback warning.
\DeclareFontShape{TU}{inconsolata(0)}{m}{it}{<->ssub*inconsolata(0)/m/n}{}

%% ---- AAMAS submission: double-blind review, ACM sigconf (acmart) ----
\settopmatter{printacmref=false}
\setcopyright{none}
\renewcommand\footnotetextcopyrightpermission[1]{}
\pagestyle{plain}

\let\Bbbk\relax
\usepackage{amsmath,amssymb,amsthm}
\usepackage{booktabs}
\usepackage{graphicx}
\usepackage{subcaption}
\usepackage{multirow}
\usepackage{array}
\usepackage{tabularx}
\usepackage{xcolor}
\usepackage{enumitem}
\usepackage{pifont}
\usepackage[capitalise,nameinlink]{cleveref}
\usepackage{caption}
\usepackage{float}
\usepackage{placeins}
\usepackage{stfloats}
\usepackage{cuted}
\usepackage{listings}
\usepackage{algorithm}
\usepackage[noend]{algpseudocode}
\lstdefinestyle{applisting}{
  basicstyle=\ttfamily\scriptsize,
  identifierstyle=\ttfamily\upshape,
  keywordstyle=\ttfamily\bfseries\upshape,
  commentstyle=\ttfamily\upshape,
  stringstyle=\ttfamily\upshape,
  breaklines=true,
  columns=fullflexible,
  keepspaces=true,
  showstringspaces=false,
  frame=single,
  rulecolor=\color{black!25},
  framerule=0.35pt,
  aboveskip=5pt,
  belowskip=5pt,
  xleftmargin=1.5pt,
  xrightmargin=1.5pt,
}
\lstdefinelanguage{Julia}{
  morekeywords={function,end,using,struct,agent,if,return,for,in,true,false,nothing},
  sensitive=true,
}
%% appendix: supplementary material separated from the eight-page body
\newcommand{\appendixfigurelayout}{%
  \clearpage
  \onecolumn
  \raggedbottom
}
\newcommand{\appendixbreak}{%
  \FloatBarrier
}

%% ---- theorem environments ----
\newtheorem{theorem}{Theorem}
\newtheorem{proposition}{Proposition}
\newtheorem{corollary}{Corollary}
\newtheorem{lemma}{Lemma}
\theoremstyle{definition}
\newtheorem{definition}{Definition}
\theoremstyle{plain}
\newtheorem*{proposition*}{Proposition}
\theoremstyle{remark}
\newtheorem{remark}{Remark}


%% number theorem-like environments as 3.<n> to match the manuscript's Section 3 numbering
\renewcommand\thetheorem{3.\arabic{theorem}}
\renewcommand\theproposition{3.\arabic{proposition}}
\renewcommand\thecorollary{3.\arabic{corollary}}
\renewcommand\thedefinition{3.\arabic{definition}}
\renewcommand\thelemma{3.\arabic{lemma}}


%% cleveref names
\crefname{theorem}{Theorem}{Theorems}
\crefname{proposition}{Proposition}{Propositions}
\crefname{corollary}{Corollary}{Corollaries}
\crefname{definition}{Definition}{Definitions}

\begin{document}

\title{AMBER: Activation Semantics for Columnar Agent-Based Modelling at Scale}

%% Replace TBD with the assigned venue tracking number before submission.
\newcommand{\submissionid}{TBD}
\author{Paper ID: \submissionid}
\renewcommand{\shortauthors}{Paper ID: \submissionid}

\begin{abstract}
Columnar and GPU execution can remove Python interpreter bottlenecks in
agent-based models (ABMs), but a performance refactor may also change the
activation regime: a snapshot update reads one step-entry state, whereas an
in-place schedule can expose earlier same-step writes. We separate four layers
that are often conflated. First, for an event multiset fixed at step entry, we
prove that snapshot execution agrees with every sequential event order when no
event reads another's target and shared targets use one associative--commutative
reducer. Second, AMBER implements a bounded runtime diagnostic for operations
observable at its API seams, including duplicate ordinary writes,
commit-before-borrow hazards, structural changes, and untraceable mutable-array
access; a clean report is not a proof about arbitrary array code. Third,
controlled ring-lattice SIR programs show that replacing the declared snapshot
rule by an in-place random-sequential rule shifts a finite-horizon attack-rate
crossing by $6.9\%$ under per-step reshuffling; a fixed-order robustness
condition and a mesa-frames blocking study show the same direction. Fourth, an
approval-gated GPU path uses AMBER's public API. On one RTX~5090, a corrected
ten-run benchmark at $10^7$ agents finds speedups of $1.77$--$2.05\times$ over
our corresponding FLAME GPU~2 implementations on three directly comparable
workloads; a setup-inclusive Schelling row is reported separately. Together
these layers connect a formal guarantee, observable reports, scientific
effects, and speed without using one as evidence for another.
\end{abstract}

\begin{CCSXML}
<ccs2012>
<concept>
<concept_id>10010147.10010919.10010177</concept_id>
<concept_desc>Computing methodologies~Multi-agent systems</concept_desc>
<concept_significance>500</concept_significance>
</concept>
<concept>
<concept_id>10010147.10010341.10010342</concept_id>
<concept_desc>Computing methodologies~Simulation theory</concept_desc>
<concept_significance>500</concept_significance>
</concept>
<concept>
<concept_id>10011007.10011074.10011099.10011692</concept_id>
<concept_desc>Software and its engineering~Formal software verification</concept_desc>
<concept_significance>300</concept_significance>
</concept>
</ccs2012>
\end{CCSXML}

\ccsdesc[500]{Computing methodologies~Multi-agent systems}
\ccsdesc[500]{Computing methodologies~Simulation theory}
\ccsdesc[300]{Software and its engineering~Formal software verification}

\keywords{agent-based simulation; execution semantics; vectorization;
runtime monitoring; emergent behaviour}

\maketitle


% ===== BODY =====

\section{Introduction}

Agent-based models (ABMs) derive collective behaviour from local rules, as in
epidemic, segregation, market, and opinion models. Large populations,
stochastic replications, and parameter sweeps make such studies expensive. In
Python, an object-level scheduler sends every agent update through the
interpreter.

Columnar CPU or GPU operations remove much of that overhead, but data layout is
only half of the translation. An ABM step also has an \emph{activation regime}.
Snapshot activation reads step-entry state; a schedule can expose same-step
writes. Replacing an immutable input buffer with an in-place
or blockwise update can therefore change the model silently. For example, a
newly infected agent may transmit immediately under one regime but only at the
next step under another.

We follow that semantic choice through four evidence layers. A reduction-aware
form of Bernstein non-interference first gives a sufficient condition for every
ordering of a \emph{fixed} event multiset to match snapshot execution. AMBER
then records hazards observable at its OOP, columnar, and device seams, while
marking mutable-array access as uncertified. Controlled SIR programs quantify
the consequence of changing activation. Finally, an independently tested GPU
path asks whether explicit semantics can coexist with population-scale
throughput. Each layer has a distinct scope.

\textbf{Contributions.}

\begin{enumerate}
\item \textbf{A sufficient condition for order independence (Section~3).}
For events fixed at step entry, reduction-aware non-interference makes snapshot
execution equal every sequential order. We extend the result to trajectories
and derive a double-buffer construction and four-class taxonomy.

\item \textbf{A bounded runtime diagnostic (Sections~3--5).} The per-step
monitor observes AMBER's buffered, columnar, and device seams. It reports
duplicate/cross-seam writes, late borrows, structural changes, and mutable-array
access. Reports concern the observed trace, not hidden indexing or the private
GPU loop.

\item \textbf{Semantic evidence (Sections~5.1--5.2).} Row-wise CPU
and batched GPU implementations of one snapshot SIR rule give consistent
crossing estimates. A matched in-place ordered variant moves the crossing by
$6.9\%$ of the snapshot estimate when the order is reshuffled every step; a
fixed order and a mesa-frames blocking study move in the
same direction. These model-level experiments are not theorem instances or
monitor runs.

\item \textbf{A corrected benchmark.} In Section~5.4, ten retained runs per
configuration at $N{=}10^7$ put AMBER GPU ahead of our FLAME GPU~2
implementations by $1.77$--$2.05\times$ on wealth transfer, random walk, and
SIR. A $63.4\times$ Schelling row includes a Python per-agent FLAME setup loop
inside the timed region and is therefore exploratory. Raw samples, medians,
IQRs, and bootstrap median intervals are released. We treat these as
implementation comparisons, not universal framework rankings.
\end{enumerate}

Figure~\ref{fig:arch} contrasts the OOP and array emphases in our released
examples and places AMBER's hybrid API between them.

\begin{figure*}[t]
\centering
\includegraphics[width=1.00\textwidth]{figs/plot01.pdf}
\caption{AMBER combines agent-centric/OOP and columnar/accelerator execution
under one model API to pursue both performance and explicit activation
alignment. The side groupings describe the emphasis of released examples, not
universal framework limits. OOP/vectorized development runs are instrumented;
private optimized GPU loops require an explicit per-instance approval label.}
\Description{A horizontal vector schematic places AMBER between agent-centric and array-or-accelerator emphases. The left box lists AgentPy, Mesa, Agents.jl, Melodie, and SimPy as released object-or-struct-per-agent examples. The central AMBER box shows one model and run API exposing OOP, vectorized-columnar, and GPU modes; OOP and vectorized modes are marked instrumented, while private GPU loops are marked approval-gated. The right box lists mesa-frames, FLAME GPU 2, and capability-only AgentTorch as released column, tensor, or device examples. A footer states that placement is not a universal capability limit.}
\label{fig:arch}
\end{figure*}

\begin{table*}[t]
\centering
\small
\caption{Capabilities represented in the released comparison (Section~5.4).
\ding{51}~= represented in the cited release/API; \ding{55}~= not identified
there and is not an impossibility claim. \emph{Best $N$} is the largest
population completed by the released implementation under the stated budget.
Contract columns describe API features, not confluence proofs.}
\label{tab:capabilities}
\setlength{\tabcolsep}{4pt}
\renewcommand{\arraystretch}{1.12}
\newcolumntype{Y}{>{\centering\arraybackslash}X}
\begin{tabularx}{\textwidth}{@{} >{\raggedright\arraybackslash}p{2.15cm} *{4}{Y} Y *{3}{Y} @{}}
\toprule
& \multicolumn{4}{c}{\textbf{Performance}} & \multicolumn{4}{c}{\textbf{Runtime checking}} \\
\cmidrule(lr){2-5}\cmidrule(lr){6-9}
\textbf{Framework} &
  {\centering Columnar / \\ vectorized\par} &
  {\centering GPU \\ backend\par} &
  {\centering Best $N$ \\ (speed)\par} &
  {\centering Calibr. \\ built in\par} &
  {\centering Diff. \\ engine\par} &
  {\centering Conflict \\ checks\par} &
  {\centering Runtime \\ monitor\par} &
  {\centering Raw-array\\warning\par} \\
\midrule
\textbf{AMBER (ours)} & \ding{51} & \ding{51} & $10^7$ & \ding{51} & \ding{55} & \textbf{\ding{51}} & \textbf{\ding{51}} & \textbf{\ding{51}} \\
FLAME GPU~2 & \ding{51} & \ding{51} & $10^7$ & \ding{55} & \ding{55} & \ding{55} & \ding{55} & \ding{55} \\
AgentTorch$^{\ddagger}$ & \ding{51} & \ding{51} & --- & \ding{51}$^{\nabla}$ & \ding{51} & \ding{55} & \ding{55} & \ding{55} \\
mesa-frames & \ding{51} & \ding{55} & $10^6$ & \ding{55} & \ding{55} & \ding{55} & \ding{55} & \ding{55} \\
Agents.jl & \ding{55} & \ding{55} & $10^6$ & \ding{55} & \ding{55} & \ding{55} & \ding{55} & \ding{55} \\
Melodie & \ding{55} & \ding{55} & $10^6$ & \ding{55} & \ding{55} & \ding{55} & \ding{55} & \ding{55} \\
AgentPy & \ding{55} & \ding{55} & $10^5$ & \ding{55} & \ding{55} & \ding{55} & \ding{55} & \ding{55} \\
SimPy & \ding{55} & \ding{55} & $10^5$ & \ding{55} & \ding{55} & \ding{55} & \ding{55} & \ding{55} \\
Mesa & \ding{55} & \ding{55} & $10^4$ & \ding{55} & \ding{55} & \ding{55} & \ding{55} & \ding{55} \\
\bottomrule
\end{tabularx}

{\footnotesize $^{\ddagger}$~AgentTorch is included for capability comparison
(differentiable GPU ABM); it was not in our matched speed run.
$^{\nabla}$~AgentTorch supports gradient-based calibration; AMBER's reported
calibration experiment uses black-box SMAC.}
\end{table*}


\section{Background and Related Work}

An ABM evolves a population state by repeatedly applying local rules. Three
literatures locate our contribution: scalable ABM execution, activation
semantics, and dependence or runtime checking. AMBER connects them, but does
not collapse them: the general runner exposes an API-level diagnostic, whereas
the optimized GPU path is assessed with model-specific reference tests rather
than a runtime certificate.

\subsection{Agent-based modelling at scale}

Several Python ABM libraries expose object-per-agent scheduler loops
\citep{kazil2020,foramitti2021,yu2023}; established toolchains use related
per-agent execution models \citep{wilensky1999,north2013}. Our timed comparison
also includes an implementation using SimPy's discrete-event processes
\citep{simpy2026}; this is a benchmark row, not a claim
that SimPy is specialized for fixed-step ABMs. Agents.jl provides a compiled
Julia implementation rather than a Python interpreter loop
\citep{datseris2022}, while Repast HPC addresses distributed-memory scale-out
\citep{collier2013}. In our particular wealth implementations and fixed time
budget, the non-GPU baselines stop between $10^4$ and $10^6$ agents; this is a
measurement of the released programs, not a general framework ceiling. A 2024
multi-engine study likewise argues for reliability tooling alongside raw speed
\citep{antelmi2024}.

Two designs change data layout and execution together.
\emph{Columnar (vectorized)} execution stores one array per attribute and
applies batched operations. Mesa-frames \citep{mesaframes} and AMBER's
vectorized CPU mode use this pattern. \emph{GPU} execution maps operations to
device kernels, continuing earlier data-parallel ABM work
\citep{lysenko2008,richmond2010}. FLAME GPU~2
\citep{brayshaw2021,richmond2023} is designed for very large GPU-resident
populations, while AgentTorch and related work
\citep{chopra2023,chopra2025,querabofarull2023} additionally support
differentiable simulation. In our RTX~5090 campaign, the four AMBER GPU
implementations are faster than our FLAME GPU~2 counterparts, but the
setup-inclusive Schelling ratio has a different evidential status from the
other three workloads (Section~5.4).

A common vectorized design evaluates a batch against shared input arrays, but
frameworks also support layering, messaging, and explicit scheduling. Those
mechanisms do not by themselves establish that a refactored program preserves
the activation regime intended by its author. Section~2.2 explains why that
distinction is consequential.

\subsection{Activation semantics and emergent behaviour}

An ABM's \emph{activation regime} specifies update order and simultaneity; it
can change the emergent outcome
\citep{axtell2001,railsback2019,grimm2006,grimm2010}. Classic digital-simulation
studies show that imposed time discretization can alter aggregate conclusions
\citep{huberman1993}. Simultaneous, uniform, and random schedules are
first-class modelling choices; segregation and cellular automata likewise show
that synchronous and sequential updates need not agree
\citep{schelling1971,fates2014}. In an epidemiological ABM, an
in-place update can let an agent infected earlier in a step transmit within
that same step, the mechanism isolated in Sections~5.1--5.3.

Snapshot-style columnar and GPU updates can implement simultaneous activation:
every event reads a common step-entry state. This property is not automatic;
an array program may still update in place or process local blocks
sequentially. Where event order is immaterial, snapshot evaluation can be
parallelized without changing the event-level result. Otherwise snapshot and
scheduled activation specify different models. Existing guidance recommends
double buffering for synchronous updates \citep{railsback2019,grimm2010}.
Section~3 gives a sufficient event-level condition and an explicit two-buffer
construction; it does not claim necessity.

\subsection{Parallel-execution conditions and runtime verification}

The formal root is Bernstein's classical parallel-execution conditions
\citep{bernstein1966}: write sets may not conflict, and no computation may read
another's writes. The same non-interference idea recurs in bulk-synchronous
semantics and runtime loop-dependence tests
\citep{valiant1990,rauchwerger1995}; deterministic-by-default parallel models
make the complementary language-design argument \citep{bocchino2009}. Dynamic
race analyses range from lockset checking through precise happens-before and
production detectors \citep{savage1997,flanagan2009,serebryany2009}. AMBER's
API monitor is narrower than those general thread-race tools. We adapt a
reduction-aware condition to one ABM event multiset fixed at step entry. Theorem~\ref{thm:char}
gives a sufficient condition under which snapshot activation equals every
scheduled event order. Appendix~\ref{ax:tax} turns the condition into a
conservative rule taxonomy, and Theorem~\ref{thm:conf} lifts the semantic
premise to a whole trajectory. The implementation monitor is narrower: it
reports selected conflicts at instrumented API seams.

That framing draws on design-by-contract \citep{meyer1992} and runtime
verification \citep{bartocci2018,leucker2009}. AMBER's diagnostic
(Appendix~\ref{ax:cert}) checks concrete API operations; Proposition
\ref{prop:errorsound} states only that reports correspond to those observed
operations. It is neither a sound-and-complete decision procedure for semantic
equivalence nor a monitor for arbitrary kernels. Its bookkeeping cost is
$O(N+q+c)$ for population size $N$, observed buffered writes $q$, and schema
columns $c$ checked at the step boundary (Section~5.3). Differentiable ABMs
\citep{chopra2023,chopra2025,querabofarull2023} address the complementary goal
of gradient-based fitting rather than this API-level diagnostic.

Among the population-scale engines in Table~\ref{tab:capabilities}, FLAME
GPU~2 \citep{richmond2023} provides compiled kernels, layering, and explicit
messaging; mesa-frames \citep{mesaframes} vectorizes Mesa through Polars; and
AgentTorch \citep{chopra2023,chopra2025} focuses on differentiable simulation.
In the cited releases and documentation, we did not identify an equivalent API
report combining duplicate ordinary writes, post-commit borrows, and
untraceable mutable access. Our mesa-frames
experiment shows that replacing a whole-array local snapshot computation with
blockwise in-place computation before the same single whole-column commit can
move an emergent threshold without an API warning. The two programs specify
different operational schedules; this is not a framework bug.

The gap is therefore not another scheduler feature. It is a disciplined link
from a formal sufficient condition, through a deliberately bounded runtime
observation, to controlled semantic evidence and an approval-gated fast path.
The monitor can expose otherwise silent hazards, but a clean trace neither
proves schedule equivalence nor certifies the private benchmark loop; the
released approval label records a deployment decision rather than a theorem.


\section{Execution semantics and the runtime diagnostic}
\label{sec:theory}

We model one simulation step as a multiset $W$ of write events generated from
the step-entry state $S^{(t)}$. Event identities, targets, and read sets are
therefore fixed before execution; rules that generate new events from same-step
writes require a separate expansion or staging argument. \emph{Scheduled}
activation applies $W$ sequentially to a running state in some order, whereas
\emph{snapshot} activation evaluates every event against $S^{(t)}$ and resolves
each target independently. Snapshot activation is one semantics that a
vectorized backend can implement, not a property of array storage itself.

A sufficient condition for agreement under every order is
\emph{non-interference}: no event reads another event's target, and all events
sharing a target use one associative--commutative (AC) reducer. This is a
reduction-aware adaptation of Bernstein's conditions (Theorem~\ref{thm:char};
full proof in Appendix~\ref{app:theory}). It is not an iff characterization:
interfering events can agree at a particular state, and non-AC binary notation
can induce commuting unary updates. Appendix~\ref{ax:tax} develops the induced
P/G/R/U rule taxonomy; the system method below operationalizes the condition
without treating that conservative taxonomy as a complete decision procedure.

For an event $w$, let $\mathrm{tgt}(w)$ be its target cell and
$\mathrm{rd}(w)$ its (possibly conservative) read set.

\begin{definition}[non-interference]\label{def:noninterference}
A fixed event multiset $W$ is non-interfering when, for every distinct
$w,w'\in W$, neither event targets a cell read by the other, and any events
sharing a target are reductions using the same associative--commutative
operator.
\end{definition}

\begin{theorem}[sufficient condition for sound vectorization]\label{thm:char}
If $W$ is fixed at $S^{(t)}$ and is non-interfering, then
\[
\mathrm{Snap}(W)(S^{(t)})=\mathrm{Sched}_{\pi}(W)(S^{(t)})
\quad\text{for every event order }\pi.
\]
\end{theorem}

\begin{proof}[Proof sketch]
Earlier events cannot alter a cell read by the current event, so every
contribution equals its value at $S^{(t)}$. Untouched cells and cells with one
target therefore agree immediately; for a shared target, the
common associative and commutative reducer makes the fold invariant to
permutation. The full definitions, proof, caveats, and counterexamples to the
converse are in Appendix~\ref{ax:char}.
\end{proof}

Non-interference is a semantic property of read and write sets. AMBER cannot
recover those sets from arbitrary Python or CUDA code, so its runtime component
checks a narrower operational contract. It records ordinary and reduction
writes routed through the public seams, flags duplicate or cross-path ordinary
writes and a borrow after a same-step commit, and reports mutable raw-array
access as uncertified. Schema/population changes are also reported. A clean
report therefore means ``no monitored violation was observed''; it does
not entail that every sequential activation order agrees. The formal whole-run
theorem is conditional on establishing non-interference at every step
(Theorem~\ref{thm:conf}), while a common synchronous gather admits a
one-barrier double-buffer construction (Theorem~\ref{thm:stage}). Monitor work is
$O(N+q+c)$ for $q$ observed buffered writes and $c$ schema columns checked
(Proposition~\ref{prop:cost}).


\begin{strip}
  \section{AMBER methodology and system realization}
  \label{sec:amber}
  Figure~\ref{fig:method} orders AMBER's evidence workflow; Algorithm~\ref{alg:amber}
  realizes its instrumented runner and gated fast path. Theorem, runtime,
  equivalence, and timing remain distinct claims.

  \centering
  \includegraphics[width=\textwidth]{figs/plot02.pdf}
  \captionsetup{hypcap=false}
  \captionof{figure}{AMBER's development-to-deployment flow. Hazards return to
    repair; external evidence supports caller approval; runtime eligibility
    selects the general or private loop. Reports, evidence, approval, and
    timings remain distinct.}
  \Description{A decision-driven flowchart. Model intent leads to selection of native OOP CPU, vectorized CPU, or GPU placement and then to an instrumented contract run. A hazard decision sends failures through repair or staging and back to the contract run. Clean runs proceed to workload evidence that may be coupled, distributional, or invariant-based. The caller either revises the model or records an approval label; a separate runtime eligibility decision selects the general native runner or the approval-gated private optimized loop. Both branches end by synchronizing, assembling the result, and retaining timing evidence.}
  \label{fig:method}
\end{strip}

\begin{algorithm}[H]
  \caption{AMBER instrumented execution and gated fast path}
  \label{alg:amber}
  \scriptsize
  \begin{algorithmic}[1]
    \Require model $M$, placement $p$, horizon $T$, contract mode $c$
    \Ensure result $R$; reports $\mathcal{C}$ (empty when $c=\texttt{off}$)
    \State $(d,m)\gets\Call{Resolve}{p}$; $\mathcal{C}\gets[\,]$; \Call{BeginExecution}{$M$}
    \If{\Call{FastEligible}{$M,d,m,c$} \textbf{and} \Call{Approved}{$M$}}
      \State $a\gets\Call{ApprovalLabel}{M}$ \Comment{non-empty, caller supplied}
      \State \Call{SetupFast}{$M,d$}
      \For{$t=0,\ldots,T-1$}
        \State \Call{FastStep}{$M,d$}
      \EndFor
    \Else
      \For{$t=0,\ldots,T-1$}
        \If{$c\neq\texttt{off}$}
          \State $C_t\gets\Call{BeginReport}{\text{schema},\text{IDs}}$
        \EndIf
        \State \Call{Dispatch}{$M,m$} \Comment{OOP or vectorized step hook}
        \If{$c\neq\texttt{off}$}
          \State $C_t\gets\Call{FinalizeReport}{C_t,\text{writes},\text{borrows},\text{drift}}$
          \State append $C_t$ to $\mathcal{C}$; \Call{ApplyPolicy}{$c,C_t$}
        \EndIf
        \State advance time; collect requested reporters
      \EndFor
    \EndIf
    \State \Call{EndExecution}{$M,d$} \Comment{synchronizes device state}
    \State \Return $(\Call{Assemble}{$M$},\mathcal{C})$
  \end{algorithmic}
\end{algorithm}

\textbf{Eligibility and policy.} Fast-path eligibility requires GPU/vectorized
placement, contract \texttt{off}, no initial record, progress, or reporters,
private setup/step hooks, and a non-empty approval label installed through
\texttt{approve\_fast\_path(evidence)}. The label is caller-supplied
provenance; \texttt{run()} checks its presence but does not validate the
evidence. Without it, the same public call uses the general runner. Reports are
stored in all instrumented modes; \texttt{warn} also emits them and
\texttt{raise} stops only on errors. Execution teardown synchronizes GPU state
before assembly.

\textbf{Native modes.} \texttt{cpu()} dispatches tracked Python \texttt{Agent}
objects in \texttt{oop} mode or Polars columns in \texttt{vectorized} mode;
\texttt{gpu()} stores numeric columns in CuPy. All share setup/result APIs and
may use mode-specific step hooks.

\textbf{Observed write seams.} OOP, vectorized, and device assignments enter a
tracked queue, frame setter, or scatter seam; incompatible cross-seam writes are
rejected. Immutable column and TensorLane reads count as borrows. Raw mutable
array access cannot be reconstructed and therefore yields an uncertified-borrow
warning. Deprecated wholesale frame replacement is outside the conflict seam,
although endpoint schema/id changes remain visible.

\textbf{Report meaning.} At step entry the monitor snapshots schema and agent
IDs. During the step it records target cells for buffered writes and column
identities for vectorized/device operations. Errors denote concrete conflicts
at those seams; warnings denote additions, births/deaths, or untraceable mutable
access. A \emph{clean} record has no errors or warnings; an \emph{OK} record has
no errors. Neither property is complete for cross-agent indexing hidden inside
user arrays. Endpoint structural checks can also miss a mutate-and-revert
sequence within one step.

\textbf{Four contract modes and a separate fast path.} The \texttt{contract}
argument accepts \texttt{off}, \texttt{check}, \texttt{warn}, and
\texttt{raise}. \texttt{check} stores reports, \texttt{warn} also emits them,
and \texttt{raise} stops on an error (warnings alone do not raise). The general
runner uses the same step hook in these modes. With GPU vectorization, contract
\texttt{off}, no initial record, no reporters or progress output, and available
private setup/step hooks, the public API selects a model-specific loop only
after explicit per-instance approval. Headline throughput therefore carries no
runtime certificate, and approval itself is not semantic validation. In this
paper the associated evidence consists of small workload smoke and invariant
tests, not a coupled-equivalence suite for the timed kernels.

Appendix~\ref{app:usage} shows all three native modes and the separate optimized
benchmark path.


\section{Experiments}

The evaluation follows the paper's evidence layers: semantic sensitivity,
monitor behavior, and throughput.
Sections~5.1--5.2 use standalone NumPy/CuPy prototypes to compare declared
snapshot updates with in-place ordered variants; they do not execute AMBER's
contract. Section~5.3 evaluates the monitor only at the API seams it observes.
Section~5.4 benchmarks the private optimized GPU path. Our corrected comparison
with FLAME uses one RTX~5090 (32\,GB, sm\_120), CUDA driver
595.71, 50 steps, one untimed warm-up, and ten retained runs per model with no
outlier deletion. We report arithmetic means in Table~\ref{tab:flame} and
release all samples, medians, IQRs, and bootstrap
median intervals, following guidance to expose variability and experimental
context in parallel-system benchmarking \citep{hoefler2015}. Other figures
retain their own stated protocols and should not be read as one pooled
benchmark campaign.

\subsection{Activation semantics move an operational epidemic crossing}

Row-wise CPU and batched GPU implementations of the same snapshot rule give
similar crossing estimates; changing to an in-place ordered rule produces a
larger, directionally consistent shift.

\begin{figure*}[t]
\centering
\includegraphics[width=\textwidth]{figs/plot03.pdf}
\caption{Ring-lattice SIR semantic sensitivity ($N{=}4000$, 120 steps).
The row-wise snapshot reference and batched CuPy snapshot implementation have
overlapping crossing intervals;
the in-place ordered implementation lets newly infected agents affect later
updates in the same step and crosses attack${=}0.5$ at about 7\% lower
transmissibility. These
are standalone experiment programs, not AMBER contract certificates.}
\Description{Three attack-rate curves over transmissibility. Row-wise CPU snapshot and batched GPU snapshot nearly overlap, while the in-place ordered curve crosses attack rate 0.5 at a lower transmissibility.}
\label{fig:threshold}
\end{figure*}

We sweep 16 transmissibility values and locate the attack-rate crossing in
three implementations, with 2000-resample bootstrap intervals. The CPU
snapshot reference reads an immutable entry copy (48 seeds); a seed-batched
CuPy version implements the same rule (96 seeds); and an ordered loop updates
in place (48 seeds).
For each CPU seed and transmissibility, the row-wise and in-place states share
the initial state, a shared activation order reshuffled at every step, and
per-agent infection/recovery draws; their contrast therefore isolates the
update regime under a random sequential schedule. Bootstrap resampling retains
each seed's trajectory across all transmissibility values. The CuPy implementation
uses an independent device RNG and serves as a distributional snapshot
cross-check. Accordingly, ``reference'' below means reference for the declared
snapshot model, not a universally correct ABM schedule.
Here, $\tau_c$ denotes the finite-$N$, 120-step operational estimate of the
attack${=}0.5$ crossing, not an asymptotic critical point.
Because same-step infections can alter later infection opportunities in the
in-place program, event eligibility is not fixed at step entry. This experiment
therefore compares two model-level activation regimes; it is not a direct
instance of Theorem~\ref{thm:char}.

\begin{table}[t]
\centering\small
\caption{Operational crossing estimate $\tau_c$ under the three execution regimes
(attack${=}0.5$ crossing; bootstrap interval over seeds).}
\label{tab:tc}
\begin{tabular}{@{}l l@{}}
\toprule
Engine & $\tau_c$ (attack${=}0.5$), median [95\% CI] \\
\midrule
Row-wise snapshot reference & 0.2422 [0.2350, 0.2509] \\
\textbf{Batched GPU snapshot} & \textbf{0.2399 [0.2348, 0.2447]} \\
Batched GPU snapshot @ $N{=}10^5$ & 0.2433 [0.2407, 0.2461] \\
In-place ordered update & 0.2257 [0.2182, 0.2330] \\
\bottomrule
\end{tabular}
\end{table}

The snapshot estimates are $0.2422$ and $0.2399$; overlapping bootstrap
intervals indicate consistency, not equivalence. The in-place estimate is
$0.2257$. In the paired CPU contrast, the absolute shift is $0.0167$
[0.0108, 0.0236], or $0.0167/0.2422=6.9\%$ of the row-wise estimate. Reusing
one fixed activation order for the full trajectory yields a larger $12.1\%$
shift [absolute interval 0.0195, 0.0390], so the direction is robust but its
magnitude is schedule-dependent. This establishes sensitivity for this finite
model and protocol, not a general claim that sequential activation is wrong. The
$N{=}10^5$ snapshot row is a scale check only; it does not justify
extrapolation to $10^7$. Crossing sensitivity at attack rates 0.3 and 0.7 is
reported in Figure~\ref{fig:sens}.

\paragraph{Released-framework sensitivity.}
Using mesa-frames' public Polars API, we implemented the same ring model with
24 seeds per transmissibility value. Each
step copies status to a local NumPy array, processes it
either as one whole-array pass or as 10, 40, 160, or 640 in-place local blocks,
and then performs one final whole-column \texttt{set}. The one-block case gives
$\tau_c{=}0.2406$; with more local blocks, later blocks see earlier local writes
and the estimate decreases monotonically to $0.2242$
(Figure~\ref{fig:anchor}). Thus every variant makes exactly one framework
commit per step; the schedule difference is in the user-side computation. The
paired-seed bootstrap intervals are [0.2251, 0.2553] and [0.2157, 0.2360] for
one and 640 blocks, respectively. This is not a mesa-frames correctness bug,
but it shows that a local refactor can change the modeled regime without an API
warning.

\subsection{Additional semantic-sensitivity prototypes}

Standalone array prototypes probe whether the same concern appears beyond the
controlled crossing experiment; their designs support only exploratory claims.

\begin{figure*}[t]
\centering
\includegraphics[width=\textwidth]{figs/plot21.pdf}
\caption{Exploratory NumPy/CuPy sensitivity studies. Snapshot/staged and
in-place variants differ in (a) the fitted SIR phase boundary, (b) Schelling
segregation, and (c) bounded-confidence consensus. Panel (c) uses unmatched
population sizes, so its gap is descriptive rather than a controlled effect.}
\Description{Three small panels compare snapshot or staged execution with in-place execution for an SIR phase boundary, Schelling segregation, and bounded-confidence consensus.}
\label{fig:sci}
\end{figure*}

The phase-boundary prototype gives slopes $-1.923$ and $-1.839$ for snapshot and
in-place variants, respectively; the mean relative crossing shift across radii
is $-7.4\%$ [bootstrap CI $-12.6\%$, $-1.5\%$]. The
Schelling prototype likewise differs near tolerance $0.75$. The coordination
prototype compares $N{=}20{,}000$ for the staged implementation with
$N{=}2{,}000$ for the in-place implementation on separately generated graphs;
finite-size and graph effects are therefore confounded with update semantics.
These studies motivate matched AMBER experiments, but the present artifacts do
not support claims that an AMBER certificate caught or prevented the effects.

\subsection{What the runtime monitor observes---and misses}

Current deterministic regression tests exercise each implemented report and
retain one explicit false-negative example. This makes the monitor's boundary
executable rather than leaving it only as a prose disclaimer.

The monitor tests executable API behavior, not the standalone science scripts
of Sections~5.1--5.2. Table~\ref{tab:detection} reports current tests against the
same repository revision as the paper. Commit-then-borrow and duplicate
ordinary writes produce errors; mutable raw-array access
produces a warning; and a single declared integer reduction remains clean. The
last row is equally important: a column borrowed once, transformed with
\texttt{np.roll}, and committed once is not reported, because the runtime does
not recover cell-level provenance from that expression. The test intentionally
asserts this clean result so future documentation cannot silently overstate the
monitor.

\textbf{Cost.} Endpoint id/schema checks cost $O(N+c)$, where $c$ is the number
of schema columns. Bookkeeping costs $O(q)$ for $q$ buffered writes. A current
$q{=}0$ microbenchmark uses one or eight commits, populations of $10^4$,
$10^5$, and $10^6$, 20 steps, and five retained runs per mode. The raw
samples and arithmetic means are released with the paper; Figure~\ref{fig:cost}
and Table~\ref{tab:scaling5} report this bounded regime. It is not directly
comparable to the private GPU loop, which disables the contract. Across these
cells, \texttt{check}/\texttt{off} ratios range from $2.95\times$ to
$18.49\times$, with 0.515--71.99 added ms per step; the monitor is therefore a
development diagnostic, not a zero-cost production feature.

\begin{table}[t]
\centering\footnotesize
\caption{Current deterministic API-level regression scenarios. ``Pass'' means
the implementation produced the documented result, including the explicit
known-limit row; it is not a completeness claim.}
\label{tab:detection}
\begin{tabularx}{\columnwidth}{@{}>{\raggedright\arraybackslash}X l l l@{}}
\toprule
Scenario & Expected & Observed & Result \\
\midrule
Commit $\rightarrow$ borrow & error & read-after-write & pass \\
Duplicate ordinary write & error & duplicate-write & pass \\
Mutable raw array & warning & uncertified & pass \\
One integer \texttt{scatter\_add} & clean & clean & pass \\
Borrow, gather, commit & no report & clean & limit \\
\bottomrule
\end{tabularx}
\end{table}

The final row is also a released end-to-end microcase. Starting from
$x=[1,2,3]$, one borrow followed by \texttt{np.roll} and one commit produces
the snapshot result $[3,1,2]$ and a clean report. Executing the analogous
predecessor rule sequentially in order $(0,1,2)$ makes all three values equal
to 3.
The external comparison therefore catches a semantic mismatch the monitor
cannot see, so no private-path approval is warranted for that translation.
This connects the workflow's report and reference roles without claiming the
monitor itself detected the dependence.

\subsection{The optimized GPU path is practical at population scale}

The full sweep places AMBER's three native modes beside seven timed baselines.
At the corrected ten-run endpoints, AMBER GPU is faster than our FLAME GPU~2
implementation on each of four $10^7$-agent workloads. The first three rows are
the directly comparable headline range; Schelling is retained as an
exploratory setup-inclusive implementation measurement.

\begin{figure*}[t]
\centering
\includegraphics[width=\textwidth]{figs/plot07.pdf}
\caption{All-framework scaling context on the RTX~5090 (50 steps). Connected
lines show the historical archived convention (mean after trimming the slowest
of ten where samples were available). A final segment connects each
series to its outlined $10^7$ AMBER GPU or FLAME GPU~2 marker, whose corrected
value is the arithmetic mean over all ten retained runs. The outlined marker
identifies the campaign/statistic change while preserving
trajectory continuity.
Missing endpoints mean no archived completion.
AgentTorch is capability-only because no matched timing row was recorded. The
Schelling label is explicitly setup-inclusive.}
\Description{Four log-log panels show historical wall-time series for wealth transfer, random walk, SIR, and Schelling across three native AMBER modes and seven external timed implementations. Corrected ten-million-agent AMBER and FLAME endpoints are outlined and connected to the preceding point; they use all-ten arithmetic means rather than the historical trimmed convention. Each panel reports the endpoint ratio, with Schelling marked setup-inclusive.}
\label{fig:gpu-benchmark}
\end{figure*}

Figure~\ref{fig:gpu-benchmark} places the corrected endpoints in the broader
scaling landscape. Table~\ref{tab:flame} reports the authoritative final means
and ratios; the released JSON retains all ten samples behind each outlined point.

\begin{table}[t]
\centering\small
\caption{Final AMBER GPU versus FLAME GPU~2 campaign at $10^7$ agents
(RTX~5090, 50 steps, ten retained runs). Entries are arithmetic means; raw
samples, medians, IQRs, and bootstrap median intervals are in the released JSON.
Speedup is FLAME/AMBER, so values above one favour AMBER. Schelling includes a
Python per-agent FLAME initialization loop in the timed region and is not part
of the headline $1.77$--$2.05\times$ range.}
\label{tab:flame}
\begin{tabular}{@{}llll@{}}
\toprule
Model & AMBER GPU & FLAME GPU~2 & speedup \\
\midrule
Wealth transfer & 94.3\,ms & 193.5\,ms & \textbf{2.05$\times$} \\
Random walk & 80.4\,ms & 160.7\,ms & \textbf{2.00$\times$} \\
SIR epidemic & 2.079\,s & 3.682\,s & \textbf{1.77$\times$} \\
Schelling (setup incl.) & 295.5\,ms & 18.72\,s & 63.4$\times$ \\
\bottomrule
\end{tabular}
\end{table}

The FLAME GPU~2 SIR row is a successful runtime-compiled execution. The archive
documents the launcher-path correction that replaced an earlier failed attempt;
the failed run is not used in Table~\ref{tab:flame}.

Timing includes construction, setup, 50 steps, and result assembly. AMBER
synchronizes before returning; FLAME returns at a completed boundary. The
timed FLAME Schelling implementation initializes agents through a Python
per-agent loop, so its $63.4\times$ ratio conflates setup strategy with step
execution and is exploratory rather than a kernel-throughput claim. The
corrected wealth kernel compacts donor eligibility before atomic debit/credit,
preventing a same-step
recipient from becoming a donor. AMBER is faster in all four measured rows
(Table~\ref{tab:flame}). These are nevertheless comparisons of our
implementations of related workloads, not identical transition kernels. The
SIR implementations use different random-number machinery, while Schelling
uses different relocation conflict handling; such choices can affect both
runtime and dynamics. AMBER's exact SIR fast path uses counting-sort cell bins
and scans candidate pairs. Its work is $O(N+C+P)$ for $C$ cells and $P$
candidate pairs examined, with a quadratic worst case if density grows in a
fixed domain; the ten-million-agent timing is therefore empirical evidence for
this configuration, not an asymptotic $O(N)$ guarantee. The FLAME models and
other baselines were implemented by one team and may not represent each
framework's maximum attainable performance.
No contract certificate is attached to these private fast-loop timings; their
released evidence is a four-workload GPU smoke/invariant test (including wealth
conservation and result shape/range checks) plus the explicit approval label.
It is not a coupled or distributional equivalence test of the timed kernels.

The ancillary studies are narrower than the original framing. Twelve KS tests
on small standalone CPU models fail to reject equality, but they are not
equivalence tests and do not cover the measured GPU kernels or Schelling. The
SLOC script counts an older public example rather than all private backend code.
SMAC improves error at matched evaluation count but takes roughly
$231$--$693\times$ more wall time than the recorded random-search runs. We retain
these artifacts in the supplement as exploratory measurements rather than
headline evidence of parity, conciseness, or practical calibration speed.

\section{Discussion and Limitations}

The results support a layered validation argument, not an end-to-end proof.
The theorem is exact but conditional on a fixed event multiset, complete read
sets, and truly AC reducers. The monitor is executable but observes only its API
seams. The SIR study identifies a causal difference between two activation
regimes under paired CPU randomness, but it does not establish that every model
is sensitive. The GPU campaign establishes implementation-level throughput on
one machine; related, rather than identical, transition kernels limit framework
ranking claims.

Monitor incompleteness is the most immediate engineering limitation. A
gather expression borrowed once before commit, value-sensitive dependence, an
unexecuted branch, mutate-and-revert structural changes, raw user kernels, and
private optimized loops can escape cell-level reconstruction. Mutable-array
warnings and device scatters avoid a vacuous clean record but do not solve
provenance. Tagged arrays or static dependence analysis could tighten
this boundary. Floating-point reductions need an orthogonal numerical policy
because machine addition is not associative.

The resulting workflow is pragmatic: declare the intended activation regime;
use \texttt{check} or \texttt{raise} while developing through observable APIs;
compare optimized paths against a coupled or distributional reference; and
report performance and semantic evidence separately. This workflow is more
informative than either a silent fast path or a certificate whose observational
scope is left implicit.

\section{Conclusion}

Activation regime is part of an ABM specification. AMBER links four scoped
forms of evidence: a sufficient event-level theorem, a bounded runtime
diagnostic, controlled semantic-sensitivity measurements, and an independently
tested GPU path. The theorem explains when fixed events are order-independent;
the monitor exposes selected development hazards; the SIR study shows their
scientific relevance; and the benchmark shows that explicit semantics can
coexist with ten-million-agent throughput. Together they make the path from
agent-centric intent to columnar or GPU execution more auditable.

\paragraph{Reproducibility.} The final AMBER--FLAME campaign uses one RTX~5090
(32\,GB, sm\_120, driver 595.71.05), Python 3.12.3, CuPy 14.1.1, a
CUDA-12.0 FLAME GPU~2 wheel (2.0.0rc4), the CUDA 12.9 host toolkit, and ten
retained runs. The released JSON contains all samples and provenance; other
studies state their own protocols.




\bibliographystyle{ACM-Reference-Format}
\bibliography{amber_references}


\appendix
\setcounter{tocdepth}{2}
\appendixfigurelayout

\section{Activation semantics: formal development and proofs}
\label{app:theory}

This appendix makes the paper's activation claim precise. It defines fixed
step-entry events and a sufficient non-interference condition, lifts the
one-step result to trajectories, and then connects the semantics to AMBER's
bounded runtime report, staging construction, and monitor cost.

\subsection{Population state and rules}\label{ax:pop}

Let the population be a set of agents $[N]=\{1,\dots,N\}$ and let $C$ be a finite set of attributes (columns), each with value domain $V_c$. A \textbf{population state} is an assignment

\[
S \in \Sigma \;:=\; \prod_{i\in[N]}\prod_{c\in C} V_c ,
\qquad S[i,c]\in V_c .
\]

A \textbf{cell} is a pair $(i,c)\in[N]\times C$. Columnar storage (a Polars DataFrame with one column per $c$) and row storage (one object per $i$) are two layouts of the same $S$; the semantics below are layout-independent.

A single \textbf{step} maps $S^{(t)}\mapsto S^{(t+1)}$. We model the step as the multiset of \textbf{write events} the rule emits when evaluated at step-entry state $S^{(t)}$. A write event $w$ carries:

\begin{itemize}
\item a \textbf{target cell} $\mathrm{tgt}(w)=(j,c)\in[N]\times C$;
\item a \textbf{mode} $\mathrm{mode}(w)\in\{\textsf{set},\
\textsf{reduce}_{\oplus_w}\}$, where $\oplus_w:V_c\times V_c\to V_c$
is the reducer carried by $w$;
\item a \textbf{read set} $\mathrm{rd}(w)\subseteq[N]\times C$ and a \textbf{value map} $\phi_w:\Sigma\to V_c$ that depends on its argument only through the cells in $\mathrm{rd}(w)$.
\end{itemize}

Read sets may be conservative supersets of the cells actually consulted. The
soundness theorem below does not need a tightness or non-triviality assumption;
we make no necessity claim. This matters because a globally non-trivial event
may be locally insensitive at $S^{(t)}$, and two interfering events can still
coincide at that state.

The rule is expanded once at step entry: fixing $S^{(t)}$ (and, for a
stochastic rule, a random tape indexed independently of execution order) fixes
the event multiset $W$, including each event's identity, target, mode, and read
set. The rule may decide which events to emit from that input, but this event
generation cannot observe same-step writes. Models with dynamic same-step event
generation fall outside the one-step theorem unless that generation is first
frozen or represented by an explicit staged semantics. Write

\[
\mathrm{WR}(W)=\{\mathrm{tgt}(w):w\in W\},\qquad
\mathrm{RD}(W)=\textstyle\bigcup_{w\in W}\mathrm{rd}(w).
\]


\subsection{Two activation semantics}\label{ax:sem}

\textbf{Snapshot (simultaneous) activation} evaluates every value map at the
frozen step-entry state $S^{(t)}$. It is a declared execution semantics, not an
automatic consequence of columnar storage. For
$W_{j,c}=\{w\in W:\mathrm{tgt}(w)=(j,c)\}$, define
$\mathrm{Snap}(W)(S^{(t)})$ cell-by-cell:

\[
S^{(t+1)}[j,c]=
\begin{cases}
S^{(t)}[j,c], & W_{j,c}=\varnothing,\\[2pt]
\phi_w(S^{(t)}), & W_{j,c}=\{w\},\ \mathrm{mode}(w)=\textsf{set},\\[2pt]
S^{(t)}[j,c]\oplus_w\phi_w(S^{(t)}),
  & W_{j,c}=\{w\},\ \mathrm{mode}(w)=\textsf{reduce}_{\oplus_w},\\[2pt]
\operatorname{fold}_{\oplus_c}\!\left(S^{(t)}[j,c],
  \{\!\{\phi_w(S^{(t)}):w\in W_{j,c}\}\!\}\right),
  & \begin{array}{l}|W_{j,c}|\ge2,\text{ all reductions use the}\\[-1pt]
    \text{same AC operator }\oplus_c,\end{array}\\[5pt]
\textbf{conflict}, & \text{otherwise.}
\end{cases}
\]

Here $\operatorname{fold}_{\oplus_c}$ starts at the entry value and combines the
multiset of contributions; AC makes its order immaterial. The conflict case
includes multiple \textsf{set}s, mixed \textsf{set}/\textsf{reduce} modes,
different reducers, and repeated non-AC reducers. The semantics deliberately
does not assign those cases an order-independent resolution.

\textbf{Scheduled (sequential) activation} --- parameterized by a total order $\pi$ on events (induced by an agent activation order). Initialize a mutable state $M\leftarrow S^{(t)}$ and process events in the order $\pi$: for event $w$, compute $v=\phi_w(M)$ against the \textbf{current} $M$ and apply it in place ---

\[
\textsf{set}:\ M[\mathrm{tgt}(w)]\leftarrow v,\qquad
\textsf{reduce}_{\oplus_w}:\ M[\mathrm{tgt}(w)]\leftarrow M[\mathrm{tgt}(w)]\oplus_w v .
\]

The output is $\mathrm{Sched}_\pi(W)(S^{(t)})=M$. Because $\phi_w$ is read against the \emph{running} state, a scheduled agent can observe earlier agents' same-step writes.

The paper's question is: \textbf{for which fixed event multisets does
$\mathrm{Snap}=\mathrm{Sched}_\pi$ for every $\pi$?} When this equality holds,
snapshot vectorization preserves the event-level result; when the equality
fails, snapshot and scheduled activation are different models.


\subsection{A sufficient non-interference theorem}\label{ax:char}

\begin{definition}[non-interference, expanded]\label{def:noninterference-full} An event multiset $W$ is \emph{non-interfering} if, for every pair of distinct events $w,w'\in W$:

\begin{itemize}
\item \textbf{(I) no cross-event read/write dependence:}
$\mathrm{tgt}(w)\notin\mathrm{rd}(w')$ and
$\mathrm{tgt}(w')\notin\mathrm{rd}(w)$; and
\item \textbf{(II) reductions only on shared targets:} if
$\mathrm{tgt}(w)=\mathrm{tgt}(w')$, then both are reductions with the
\textbf{same} operator $\oplus_c$, and $\oplus_c$ is associative and
commutative (AC).
\end{itemize}

Condition (I) is the flow/anti-dependence clause; condition (II) is the output-dependence clause relaxed to admit AC reductions. Together they are Bernstein's conditions with a reduction carve-out.

\end{definition}

\begin{proposition}[AC folds are order-invariant]\label{prop:acfold}
If all events on a cell $(j,c)$ are $\textsf{reduce}_{\oplus_c}$ with $\oplus_c$ AC and each contributes the fixed value $\delta_w=\phi_w(S^{(t)})$, the cell's final value $S^{(t)}[j,c]\oplus_c\bigoplus_w\delta_w$ is independent of the order in which the events are applied.
\end{proposition}

\begin{proof} Let $\delta_1,\dots,\delta_m$ be the increments. The scheduled fold computes $(((s\oplus\delta_{\sigma(1)})\oplus\delta_{\sigma(2)})\oplus\cdots)$ for the applied permutation $\sigma$, where $s=S^{(t)}[j,c]$. Associativity makes any re-bracketing of an iterated $\oplus$ equal, so the value depends only on the sequence $(\delta_{\sigma(1)},\dots,\delta_{\sigma(m)})$ and not on the grouping. Commutativity makes it invariant under permutation of that sequence: for any transposition of adjacent elements, $a\oplus b=b\oplus a$, and adjacent transpositions generate the symmetric group, so every $\sigma$ yields the same value. Hence the result is a function of the \emph{multiset} $\{\delta_w\}$ and $s$ alone. \end{proof}
\begin{theorem}[expanded statement of Theorem~\ref{thm:char}]\label{thm:char-full}
Let $W$ be the fixed event multiset emitted at $S^{(t)}$. If $W$ is
non-interfering, then the snapshot result is well defined and
\[
\mathrm{Snap}(W)(S^{(t)})=
\mathrm{Sched}_\pi(W)(S^{(t)})
\qquad\text{for every event order }\pi .
\]
\end{theorem}

\begin{proof} Assume (I)--(II). Before $w$ is evaluated, $w$ itself has not yet
executed, and by (I) no earlier \emph{other} event can have modified a cell in
$\mathrm{rd}(w)$. Thus $M$ and $S^{(t)}$ agree on $\mathrm{rd}(w)$, so
$\phi_w(M)=\phi_w(S^{(t)})$ independently of $\pi$. Fix a cell $(j,c)$. It is
either untouched, targeted by one event, or targeted by multiple events. The
first two cases are immediate for either mode. In the third case, condition
(II) requires one shared AC reducer, so Proposition~\ref{prop:acfold} gives an
order-invariant result and excludes the snapshot conflict case. Hence every
cell agrees under snapshot and scheduled execution for every $\pi$.
\end{proof}

\begin{remark}[why the converse is not claimed]
Interference is evidence of a potential schedule dependence, not a proof of one
at a fixed state. For example, let one event change $x:0\mapsto1$ and another
write $y:=\mathbf 1[x=2]$. The second event genuinely depends on $x$ over the
domain, yet both orders agree at this state. Likewise, two set events may write
the same value, and Proposition~\ref{prop:commuting} gives non-AC updates whose
induced unary maps commute. These counterexamples invalidate an iff statement.
\end{remark}
 
\subsection{A taxonomy of agent rules}\label{ax:tax}

Theorem~\ref{thm:char} induces a conservative classification: satisfying a
class condition is sufficient, while falling outside it calls for staging or
model-specific analysis.

\textbf{Class P --- pointwise self-maps.} Each agent $i$ emits one event with $\mathrm{tgt}=(i,c)$ and $\mathrm{rd}\subseteq\{i\}\times C$ (reads and writes only its own row). Distinct agents touch distinct target cells and read only their own, so (I)--(II) hold trivially. Such a rule satisfies the sufficient condition (e.g. per-agent age increment or independent noise on \texttt{x}).

\textbf{Class G --- snapshot gathers.} Agent $i$ writes its own cell but reads
a neighbourhood. Writing a fresh output column satisfies (I). Reading and
writing the same column fails the sufficient test and may be schedule-sensitive;
double buffering expresses the intended synchronous rule explicitly.

\textbf{Class R --- commutative reductions (many-to-one).} Several source agents contribute to one target cell through an AC reducer: resource flows, contact counts, force accumulation. Written with \texttt{scatter\_add} ($\oplus=+$), duplicate target ids satisfy condition (II); when contributions are fixed from the entry state so that (I) also holds, Theorem~\ref{thm:char} makes the update safe. This is precisely why \texttt{agents.\allowbreak{}at[\allowbreak{}[\allowbreak{}1,1,3]\allowbreak{}]\allowbreak{}.\allowbreak{}scatter\_\allowbreak{}add(\allowbreak{}wealth=\allowbreak{}1)\allowbreak{}} must credit agent 1 twice rather than clobber. \textbf{Potentially unsafe variant:} expressing the same many-to-one flow with ordinary assignment (\texttt{agents.\allowbreak{}at[\allowbreak{}ids]\allowbreak{}.\allowbreak{}wealth =\allowbreak{} .\allowbreak{}.\allowbreak{}.\allowbreak{}} with repeated ids) is a Case-II conflict. Last-writer-wins is order-dependent when repeated targets receive different values; equal-value duplicates are still rejected conservatively.

\textbf{Machine-arithmetic caveat.} The theorem requires an actually
associative and commutative operator. Fixed-width integer addition is AC on a
specified wraparound domain (for example, modulo $2^k$); overflow-raising or
saturating arithmetic requires separate analysis, and IEEE floating-point
addition is not associative. A
floating \texttt{scatter\_add} can therefore vary in low-order bits with atomic
order and requires a numerical tolerance or reproducible reduction; the API
name alone does not discharge the theorem's premise.

\textbf{Class U --- potentially order-sensitive rules.} Same-step cross-agent
read/write overlap or non-AC many-to-one writes fail the sufficient condition.
They should be staged, given an explicit schedule, or justified separately. The
runtime monitor detects selected operational manifestations, not membership in
U for arbitrary Python/CUDA code.

\begin{remark}[the AC requirement is conservative]\label{rem:ac}
Repeated subtraction induces commuting translations $x\mapsto x-\delta$ even
though the binary operator is not AC. Proposition~\ref{prop:commuting}
characterizes this family. AMBER exposes only declared AC scatter reducers, so
the API may reject such an order-safe update without implying that all non-AC
folds are unsafe.
\end{remark}

\begin{table}[H]
\centering
\small
\setlength{\tabcolsep}{4pt}
\begin{tabularx}{\textwidth}{@{}l l l >{\raggedright\arraybackslash}X >{\raggedright\arraybackslash}X@{}}
\toprule
Class & Reads & Writes & Safe to vectorize? & AMBER idiom \\
\midrule
\textbf{P} pointwise self-map & own row & own cell & always & \texttt{agents.\allowbreak{}col =\allowbreak{} f(\allowbreak{}agents.\allowbreak{}col)\allowbreak{}} \\
\textbf{G} snapshot gather & neighbourhood & own cell & safe with a fresh output column & double-buffer: read old col, write new col \\
\textbf{R} commutative reduction & non-interfering snapshot reads & shared cell, AC & safe under (I) with declared reducer & \texttt{agents.\allowbreak{}at[\allowbreak{}ids]\allowbreak{}.\allowbreak{}scatter\_\allowbreak{}add(\allowbreak{}col=\allowbreak{}delta)\allowbreak{}} \\
\textbf{U} potential conflict & --- & --- & analyze, stage, or schedule & selected hazards reported \\
\bottomrule
\end{tabularx}
\caption{The four rule classes, their read/write shapes, and the AMBER idiom for each.}\label{tab:classes}
\end{table}




\subsection{Commuting non-AC reducers}\label{ax:propA1}

Remark~\ref{rem:ac} observes that the associative--commutative requirement on
Class-R reducers is sufficient but conservative. We formalize this.

\begin{proposition}[commuting induced updates]\label{prop:commuting}
Let $D$ be a set of admissible fixed increments for reduce events on one cell,
all using a common binary operator $\oplus$, and define
$T_\delta(x)=x\oplus\delta$. For a finite sequence
$\delta_1,\dots,\delta_m\in D$, consider the scheduled left-fold
$g_\sigma(s)=T_{\delta_{\sigma(m)}}\circ\cdots\circ
T_{\delta_{\sigma(1)}}(s)$. The increments are evaluated at step entry and are
independent of the other events by clause (I).

\begin{enumerate}
\item If $\oplus$ is AC then $g_\sigma$ is independent of $\sigma$ for all
states and increments (Proposition~\ref{prop:acfold}).
\item More generally, every finite fold over increments from $D$ is
order-independent for every starting state iff the induced unary maps
$\{T_\delta:\delta\in D\}$ pairwise commute under composition. Subtraction is
the canonical example:
$T_\delta(x)=x-\delta$ are commuting translations although the binary operator
$-$ is neither associative nor commutative.
\end{enumerate}
\end{proposition}

\begin{proof} Part 1 is Proposition~\ref{prop:acfold}. For part 2,
sufficiency follows because any permutation is reached by adjacent
transpositions, each of which preserves the composition when the corresponding
maps commute. For necessity, choose any $\delta_1,\delta_2\in D$ and the
two-increment sequence. Order independence for every $s$ says exactly
$T_{\delta_1}\circ T_{\delta_2}=T_{\delta_2}\circ T_{\delta_1}$ as functions.
Thus every pair commutes. For subtraction,
$T_\delta(x)=x-\delta$ and
$T_{\delta_1}(T_{\delta_2}(x))=x-\delta_2-\delta_1=
T_{\delta_2}(T_{\delta_1}(x))$. The family therefore commutes even though the
binary operator $-$ is neither associative nor commutative, so
Definition~\ref{def:noninterference}(II) conservatively rejects it. \end{proof}

AMBER's public scatter API accepts declared AC reductions. A non-AC operation
therefore falls outside that idiom even when Proposition~\ref{prop:commuting}
shows its induced updates commute. This is an API restriction, not evidence that
the operation is schedule-dependent.

\subsection{From theorem to an operational runtime report}\label{ax:cert}

For arbitrary Python rule bodies, exact read/write provenance is unavailable at
the ordinary array API. AMBER therefore monitors concrete operations rather
than claiming to decide Theorem~\ref{thm:char}. It observes buffered OOP writes,
columnar/TensorLane commits, device scatter writes, declared reductions, and
immutable borrows. A \textbf{\texttt{ContractCertificate}} records:

\begin{itemize}
\item \textbf{duplicate ordinary write:} a buffered cell or whole column is
written more than once without using a declared reduction;
\item \textbf{cross-path write:} the same column receives incompatible writes
through different observed seams in one step;
\item \textbf{commit-before-borrow:} a column is borrowed after a same-step
commit, so that borrow cannot be guaranteed to observe entry state;
\item \textbf{uncertified mutable borrow:} \texttt{agents.array(...)} exposes a
mutable NumPy/CuPy buffer whose later indexing/mutation is not traceable;
\item \textbf{schema / population mutation:} the column set, a column dtype,
or the id-set differs between step entry and exit. Such a step is outside the
fixed-$\Sigma$ model of \S\ref{ax:pop}. Additions and births/deaths are warnings;
column removal and dtype change are errors. Because this check compares only
endpoints, a mutate-and-revert sequence can remain unobserved.
\end{itemize}

\begin{proposition}[operational meaning of reports]\label{prop:errorsound}
Each duplicate-write, cross-path-write, or commit-before-borrow error
corresponds to the concrete instrumented operations named by the report. While
monitoring is active, mutable raw-array access through the observed accessor
prevents a clean certificate. No converse or schedule-divergence conclusion is
implied without model-specific read-set information.
\end{proposition}

\begin{proof} The monitor adds a duplicate report only when the corresponding
per-step counter exceeds one, a cross-path report only when incompatible seam
records share a column, and a commit-before-borrow report only when the column
is already in the committed set. Mutable array access calls the dedicated
warning seam. These are direct invariants of the implementation. The monitor
does not inspect value sensitivity or arbitrary array indices, so stronger
semantic conclusions would not follow.\end{proof}

\begin{remark}[explicit incompleteness]
A normal gather may borrow a column once and commit once. For example, from
$x=(1,2,3)$ the snapshot assignment $x\leftarrow\mathrm{roll}(x,1)$ produces
$(3,1,2)$, while forward and reverse per-cell schedules produce $(3,3,3)$ and
$(2,1,2)$. Borrow-before-commit order alone does not reveal that cross-agent
dependence. Through \texttt{agents.array}, AMBER now reports an uncertified
mutable borrow; through an immutable but provenance-free expression, static or
tagged-array analysis would still be required.
\end{remark}
 
\subsection{Worked examples}\label{ax:examples}

\textbf{Safe construction (Class R) --- snapshot wealth transfer.} Freeze the
donor set from step-entry wealth, emit one debit per donor, and aggregate
recipient credits with addition. Fixed eligibility plus AC credit reduction
satisfies the theorem; merely checking live wealth inside concurrent atomics
would not.

\textbf{Potential conflict (Class U $\rightarrow$ fix) --- SIR.} Reading and
writing \texttt{status} in place fails condition (I). Computing
\texttt{next\_status} from an immutable entry column and swapping explicitly
implements synchronous semantics. The monitor catches commit-then-borrow
versions; raw or borrow-once kernels require separate validation.

\textbf{Potential conflict (Class U) --- duplicate ordinary assignment.} \texttt{agents.\allowbreak{}at[\allowbreak{}[\allowbreak{}1,1,3]\allowbreak{}]\allowbreak{}.\allowbreak{}wealth =\allowbreak{} v} issues two \textsf{set}s to cell $(1,\text{wealth})$: Case II. Ordinary duplicate assignment has no declared schedule-independent resolution and is potentially order-dependent when values differ; equal values are a conservative rejection. The monitor raises duplicate-unreduced-write. The sanctioned accumulation rewrite is \texttt{scatter\_add}, which is Class R when condition (I) and the arithmetic caveat also hold.


\subsection{Trajectory confluence}\label{ax:conf}

Theorem~\ref{thm:char} is a statement about one step. It lifts to trajectories
when its semantic premise---not merely a clean runtime certificate---holds at
every reached state.

Fix an initial state $S^{(0)}$, a rule-expansion map $R$, and a random tape
$\omega=(\omega_0,\omega_1,\ldots)$ indexed by step and stable event/agent
identity rather than execution position. At step $t$, expand the rule exactly
once as $W_t=R(S^{(t)},\omega_t)$. A \textbf{scheduled run} chooses an order
$\pi_t$ of $W_t$; a \textbf{snapshot run} resolves the same $W_t$ at once:
\[
S^{(t+1)}=\mathrm{Sched}_{\pi_t}(W_t)(S^{(t)})
\quad\text{vs.}\quad
S^{(t+1)}=\mathrm{Snap}(W_t)(S^{(t)}).
\]
Write a schedule as the sequence $\Pi=(\pi_0,\pi_1,\dots)$.

\begin{definition}[non-interfering run]\label{def:noninterfering-run}
A snapshot run of length $T$ is \emph{non-interfering} if, at every reached
state $S^{(t)}$, the fixed event multiset $R(S^{(t)})$ satisfies
Definition~\ref{def:noninterference}, where $R(S^{(t)})$ abbreviates
$R(S^{(t)},\omega_t)$. Thus event generation is a deterministic function of
state and the coupled tape, not activation order.
\end{definition}

\begin{theorem}[confluence]\label{thm:conf}
If the snapshot run of $R$ from $S^{(0)}$ is non-interfering for $T$ steps,
then for every schedule $\Pi=(\pi_0,\dots,\pi_{T-1})$ the scheduled run
coincides with the snapshot run state-for-state: \[
S^{(t)}_{\Pi}=S^{(t)}_{\mathrm{Snap}}\qquad\text{for all } 0\le t\le T .
\] In particular all schedules agree with each other; for stochastic rules the
trajectory is a function of $(S^{(0)},R,\omega)$.

\end{theorem}

\begin{proof} Induct on $t$. Equal states and the same $\omega_t$ emit the same
$W_t=R(S^{(t)},\omega_t)$. By
Definition~\ref{def:noninterfering-run}, $W_t$ is non-interfering, so
Theorem~\ref{thm:char} equates the snapshot update with the selected
$\pi_t$. Hence the next states agree.\end{proof}
 
\subsection{A one-barrier double-buffer construction}\label{ax:stage}

The common same-column synchronous gather can be implemented by double
buffering: read the old column, write a fresh workspace column, then commit or
swap that workspace.

Model a rule as a set of events whose reads and writes range over columns $\mathrm{Cols}=\{c_1,\dots,c_m\}$. A \textbf{staging} with $b$ barriers partitions the step into $b{+}1$ ordered phases $P_0,\dots,P_b$; each phase is executed snapshot-style on the state left by the previous phase, and a \textbf{barrier} is a point at which all writes so far are committed and become visible to later reads. A staging is \textbf{valid} if within every phase the events are non-interfering (so each phase is itself vectorizable by Theorem~\ref{thm:char}) and the composed result equals the intended \textbf{synchronous snapshot} semantics of the rule (every agent reads the step-entry neighbourhood and writes its post-step value).

Consider the canonical rule
\[
c[i]\ \leftarrow\ f\big(\,c[j] : j\in \mathcal N(i)\,\big)
\qquad\text{(read and write the same column $c$).}
\]

\begin{theorem}[one-barrier sufficient construction]\label{thm:stage}
Let $c'$ be temporary workspace outside the observable model state. The
two-phase staging $P_0=\{$read $c$, write $c'\}$, barrier,
$P_1=\{$copy $c'\to c\}$ reproduces the declared synchronous snapshot update
on the original columns and uses one inter-phase barrier.

\end{theorem}

\begin{proof} Phase 0 reads $c$ and writes only $c'$, so (I) holds; each
$(i,c')$ has one writer, so (II) holds in the state extended by the workspace.
Theorem~\ref{thm:char} yields
$c'[i]=f(c[j]:j\in\mathcal N(i))$ evaluated at step entry. Phase 1 is a
pointwise copy from $c'$ to $c$ and is likewise non-interfering. Their
composition, projected onto the original columns after discarding $c'$, is
exactly the declared snapshot update.\end{proof}

The theorem does not claim one barrier is necessary. A zero-barrier
implementation can be correct if it reads an immutable entry buffer and writes
a separate output buffer inside one kernel launch. Likewise, longest-path
layering of a dependency DAG implements topological staged semantics, not a
global step-entry snapshot; the topological-staging harness is therefore
retained only as an experiment about that different semantics.
A pointer swap between the entry and output buffers is observationally
equivalent to the explicit copy when buffer identity is not part of model state.

\subsection{A cost model for the monitor}\label{ax:cost}

The monitor combines population/schema checks with bookkeeping proportional to
the operations it observes.

Let $N$ be population size, $c$ the number of schema columns checked at the
step boundary, and $q$ the number of observed buffered cell writes. Building
and comparing endpoint id sets and schema/dtype maps costs $O(N+c)$; tracking
borrowed and committed columns costs at most $O(c)$; and the per-cell write map
costs $O(q)$ expected time under standard hash-table assumptions.

\begin{proposition}[monitor cost]\label{prop:cost}
The per-step bookkeeping work and auxiliary space of
\textnormal{\texttt{check}} mode are, respectively,
\[
O(N+q+c)\quad\text{and}\quad O(N+q+c).
\]
It does not evaluate the rule's arithmetic, but it depends on the number and
kind of observed writes. For whole-column vectorized code $q$ can be zero; for
an OOP step $q$ may be $\Theta(Nc)$.

\end{proposition}

\begin{proof} Endpoint comparisons store and examine at most $N$ ids and $c$
schema entries. Borrow/commit bookkeeping stores at most $c$ column names, with
expected constant-time hash operations. Each buffered write performs one
expected constant-time hash-map update and can add one entry, giving $O(q)$
time and space. Summing the terms proves both bounds.\end{proof}


\section{Bounded executable checks}
\label{app:referee}

The released \texttt{theorem\_referee.py} implements \texttt{Snap} and
\texttt{Sched} directly. It is useful as a bounded regression suite, but it is
not a proof assistant and does not model the implementation monitor.

\subsection{Method}\label{ax:ref:method}

The harness uses finite value domains and small populations. For each generated
event multiset it compares snapshot execution with scheduled permutations over
sampled states; selected small cases enumerate the finite state/order grid.

Three search modes are used:

\begin{itemize}
\item \textbf{Randomized states and schedules} for generated event multisets;
\item \textbf{finite exhaustive grids} where the state/order product is small;
\item \textbf{multi-step checks} for rules generated to satisfy
Definition~\ref{def:noninterference} at each reached state.
\end{itemize}

\subsection{Findings}\label{ax:ref:findings}

Table~\ref{tab:referee} reports the historical run with corrected
interpretation. None of 1327 generated non-interfering multisets violated the
sufficient direction on the tested states and orders. Among 6673 interfering
multisets, 34 dedicated value-preserving controls had no divergence on their
exhaustively enumerated finite state/order grids. Those controls demonstrate
that interference is not sufficient for divergence in the constructed cases;
they do not establish a general necessity characterization.

\begin{table}[!ht]
\centering\small
\caption{Bounded semantic checks. Counts concern the mathematical event model,
not certificates emitted by the implementation monitor.}
\label{tab:referee}
\setlength{\tabcolsep}{6pt}
\begin{tabularx}{\textwidth}{@{}l l >{\raggedright\arraybackslash}X >{\raggedright\arraybackslash}X@{}}
\toprule
Check & Cases & Finding & Interpretation \\
\midrule
Thm.~\ref{thm:char} & 1327 non-interfering & no bounded counterexample & regression-consistent with sufficient direction \\
Converse controls & 6673 interfering & 34 exhaustive finite-grid no-witness cases & constructed interference can remain order-invariant \\
Thm.~\ref{thm:conf} & 6000 premise-satisfying runs & no bounded counterexample & generated deterministic trajectories \\
Thm.~\ref{thm:stage} & ring, $N{=}6$, all $6!$ orders & double buffer agrees; zero-buffer differs & validates the construction on this grid \\
Prop.~\ref{prop:commuting} & subtraction family & no order difference & non-AC notation, commuting induced maps \\
\bottomrule
\end{tabularx}

\end{table}

\subsection{Interpretation}\label{app:ref:interp}

The development harness exposed three useful lessons:

\begin{enumerate}
\item \textbf{Understated read sets} can invalidate a theorem application even
when the runtime trace looks benign.
\item \textbf{Snapshot conflict handling} must distinguish one reduction from
mixed or repeated incompatible writes.
\item \textbf{Interference need not imply divergence.} Value-preserving writes,
local insensitivity, equal set values, and commuting non-AC maps all defeat a
general converse. The paper now states only the valid sufficient direction.
\end{enumerate}

\subsection{Scope of the evidence}\label{ax:ref:scope}

This evidence is bounded: finite values, small populations, and randomized cases
above the exhaustive grid. It is a regression check, not a proof of the general
theorem or a validation of every runtime-monitor path.

\subsection{Contrasting topological-staging experiment}\label{app:ref:barrier}

The companion \texttt{topological\_staging\_experiment.py} generates 90
dependency DAGs and compares layered execution with a \emph{sequential
topological} reference for
$x_c\leftarrow1+\sum_{p\to c}x_p$.

Longest-path layering matches that topological reference on all generated
graphs, while one fewer layer differs. Later phases intentionally read earlier
writes, so this validates topological staged semantics, not the global
step-entry snapshot semantics of Theorem~\ref{thm:stage}.

\begin{figure}[!ht]
\centering
\includegraphics[width=\textwidth]{figs/plot11.pdf}
\caption{Topological-staging experiment over 90 DAGs. The result concerns a
sequential topological reference and is not a lower bound for global snapshot
semantics.}
\Description{Bars summarize 90 dependency DAG checks. Longest-path layering matches the sequential topological reference, whereas one fewer layer can differ.}
\label{fig:barrier}
\end{figure}

\section{Supplementary experimental evidence}
\label{app:suppexp}

This appendix provides the supporting plots and tables cited by the main
experiments. Each item inherits the scope and caveats of its corresponding
section. Appendix~\ref{app:usage} then gives side-by-side wealth-transfer API
examples for the Section~5.4 implementations.

\subsection{Supporting figures for \S5.1 (threshold experiment)}\label{app:supp:sec51}

\begin{figure}[H]
\centering
\includegraphics[width=\textwidth]{figs/plot17.pdf}
\caption{mesa-frames sensitivity study. \textbf{(a)} Attack-rate curves change
as the number of local update blocks per step increases. \textbf{(b)} Recovered
$\tau_c$ drifts from the one-local-block result. Every variant performs one
final whole-column \texttt{set}; only the local NumPy computation differs.}
\Description{Left panel plots attack-rate curves for one through 640 local update blocks. Right panel plots the recovered epidemic crossing estimate against local update-block count; all variants make one final framework commit per step.}
\label{fig:anchor}
\end{figure}

\begin{figure}[H]
\centering
\includegraphics[width=\textwidth]{figs/plot18.pdf}
\caption{Crossing-level sensitivity. At attack-rate crossings 0.3/0.5/0.7,
the batched GPU snapshot implementation tracks the row-wise snapshot reference;
the in-place ordered variant is lower. At the 0.7 crossing, 1789/2000 row-wise
and 1890/2000 GPU bootstrap resamples cross within the swept range, so those
intervals are conditional on valid crossings.}
\Description{Threshold estimates with uncertainty intervals at attack-rate crossing levels 0.3, 0.5, and 0.7 for row-wise snapshot, batched GPU snapshot, and in-place ordered execution.}
\label{fig:sens}
\end{figure}

\appendixbreak
\subsection{Supporting figures for \S5.3--5.4}\label{app:supp:sec54}

The remaining figures separate four questions: CPU distributional agreement,
historical public-example size, calibration quality and cost, and runtime
monitor overhead. They are complementary diagnostics rather than a composite
framework ranking.

\begin{figure}[H]
\centering
\includegraphics[width=\textwidth]{figs/plot13.pdf}
\caption{Small-model CPU distribution checks. Three ECDF panels compare AMBER
vectorized and loop paths with AgentPy, Mesa, and mesa-frames. Across the 12
pairwise KS checks, 12/12 have $p>0.05$; failure to reject is not equivalence,
and this check does not cover the headline GPU kernels or Schelling.}
\Description{Three empirical cumulative-distribution panels compare five CPU implementations for wealth transfer, random walk, and well-mixed SIR; each panel reports the minimum pairwise KS p-value.}
\label{fig:parity}
\end{figure}

\begin{figure}[H]
\centering
\includegraphics[width=\textwidth]{figs/plot14.pdf}
\caption{Historical logical SLOC for released public benchmark examples by
framework and model. A dash denotes no implementation in the recorded
artifact. Private optimized backend code is excluded, so this is not a full
implementation-size or conciseness comparison.}
\Description{A matrix reports historical logical source-line counts for public wealth, random-walk, SIR, and Schelling examples across AMBER and seven comparison frameworks.}
\label{fig:usability}
\end{figure}

\begin{figure}[H]
\centering
\includegraphics[width=\textwidth]{figs/plot15.pdf}
\caption{Calibration at matched forward-model evaluation counts (four trials
per point). Panel (a) reports recovered-parameter $L_2$ error and trial SD;
panel (b) reports wall-clock cost. Batched SMAC records lower error at three of
four budgets but costs roughly $231$--$693\times$ more wall time than random
search.}
\Description{Two panels compare batched SMAC and random search over four matched evaluation budgets: mean parameter error with standard deviations and mean logarithmic wall time.}
\label{fig:calib}
\end{figure}

\begin{figure}[H]
\centering
\includegraphics[width=0.94\textwidth]{figs/plot12.pdf}
\caption{Fresh current-code monitor microbenchmark in the $q{=}0$ whole-column
regime. Panel (a) shows the median paired added cost and IQR over five retained
runs; panel (b) shows the arithmetic-mean wall-time ratio. The general
implementation bound remains $O(N+q+c)$ (Proposition~\ref{prop:cost}).}
\Description{Two panels show monitor overhead for one and eight columns across ten thousand, one hundred thousand, and one million agents: added milliseconds per step and check-to-off wall-time ratio.}
\label{fig:cost}
\end{figure}

\begin{table}[H]
\centering\small
\caption{Current monitor microbenchmark (20 steps, five retained runs).
Entries are arithmetic-mean \texttt{check}/\texttt{off} ratios with mean added
milliseconds per step in parentheses. Timed scope includes model construction,
setup, the step loop, and result assembly. These values do not compare against
the private headline GPU loop, which \texttt{check} disables.}
\label{tab:scaling5}
\begin{tabular}{@{}lrrr@{}}
\toprule
Whole-column commits & $10^4$ & $10^5$ & $10^6$ \\
\midrule
1 column & 10.18$\times$ (0.515) & 18.49$\times$ (6.695) & 15.38$\times$ (68.13) \\
8 columns & 2.95$\times$ (0.849) & 3.61$\times$ (7.332) & 3.68$\times$ (71.99) \\
\bottomrule
\end{tabular}
\end{table}

\appendixbreak
\subsection{Synthesis table}\label{app:supp:multiaxis}

\begin{table}[H]
\centering\footnotesize
\caption{Scope of the released artifacts. Dashes mean the corresponding
experiment did not test that path; they are not negative framework judgments.}
\label{tab:multiaxis}
\begin{tabularx}{\textwidth}{>{\raggedright\arraybackslash}X>{\raggedright\arraybackslash}X>{\raggedright\arraybackslash}X>{\raggedright\arraybackslash}X>{\raggedright\arraybackslash}X>{\raggedright\arraybackslash}X>{\raggedright\arraybackslash}X}
\toprule
Path & Wealth speed & Runtime evidence & Distribution study & Calibration & Monitor study & SLOC scope \\
\midrule
\textbf{AMBER private GPU} & 94.3\,ms & smoke/invariant tests + approval label & not tested & separate artifact & not monitored & private code excluded \\
AMBER general vectorized & 15.2\,s & API diagnostics & CPU only; no equivalence margin & separate artifact & targeted tests + current cost run & public example \\
FLAME GPU~2 & 193.5\,ms & transition tests only & not tested & --- & --- & RTC host/device split \\
mesa-frames & 23.6\,s (1M) & semantic prototype & reference only & --- & --- & public example \\
AgentPy & 122\,s (1M) & reference run & reference only & --- & --- & public example \\
Mesa & 24.9\,s ($10^4$) & reference run & reference only & --- & --- & public example \\
\bottomrule
\end{tabularx}
\end{table}

\appendixbreak
\section{Framework usage examples}
\label{app:usage}

This appendix shows the APIs used by the wealth-transfer implementations. The
benchmark's AMBER OOP, vectorized CPU, and GPU paths freeze donor eligibility
at step entry; the snapshot-safe fused CUDA kernel receives that donor list
explicitly.
Listings for external frameworks illustrate their native APIs, but transition
details (self-transfer handling, RNG, layering) are not byte-identical. Full
runnable sources, rather than these abbreviated listings, define the benchmark.

\subsection{Execution paradigms and dispatch}
\label{app:usage:overview}

AMBER exposes three native execution paths behind one public run API:
\texttt{cpu(mode="oop")}, \texttt{cpu(mode="vectorized")}, and
\texttt{gpu()}. A model may implement \texttt{step\_oop} and
\texttt{step\_vectorized} separately. For the headline GPU benchmark, the same
public \texttt{gpu().run()} call selects a private model-specific loop when the
contract and reporters are off \emph{and} the instance carries a non-empty
\texttt{approve\_\allowbreak fast\_\allowbreak path(\allowbreak evidence)}
label. The benchmark uses the released
GPU smoke/invariant suite as that label's provenance. Thus API continuity and
approval should not be mistaken for an identical step body or a runtime
certificate.

\begin{table}[H]
\centering\small
\caption{Execution paradigms represented in the wealth benchmark. AMBER shares
the public model/run API across native paths but uses path-specific hooks.}
\label{tab:usage_paradigm}
\begin{tabularx}{\textwidth}{@{} l X X @{}}
\toprule
Framework & Population representation & Step dispatch \\
\midrule
AgentPy, Mesa, Agents.jl & one object (or struct) per agent & per-agent \texttt{step} in a scheduler loop \\
mesa-frames \citep{mesaframes} & Polars \texttt{DataFrame} via \texttt{AgentSetPolars} & vectorized column \texttt{set} after NumPy work \\
AMBER (OOP) & tracked Python agent objects & \texttt{cpu(mode="oop").run()} $\rightarrow$ \texttt{step\_oop} \\
AMBER (vectorized) & Polars \texttt{DataFrame} + view API & \texttt{cpu(mode="vectorized").run()} $\rightarrow$ \texttt{step\_vectorized} \\
AMBER (GPU) & CuPy-resident numeric columns & approval + \texttt{gpu().run()} $\rightarrow$ optional model-specific loop \\
FLAME GPU~2 \citep{richmond2023} & \texttt{AgentVector} on device & layered RTC agent functions + messaging \\
\bottomrule
\end{tabularx}
\end{table}

\appendixbreak
\subsection{AMBER: OOP, vectorized CPU, and GPU}
\label{app:usage:amber}

\paragraph{OOP CPU path (\S5.4 ``loop'').}
Listing~\ref{lst:amber-loop} sketches the corrected OOP lane. Donor eligibility
and recipient draws are fixed before writes; deltas are aggregated so each
agent receives one final buffered assignment.

\begin{lstlisting}[style=applisting,language=Python,caption={AMBER wealth transfer --- OOP loop baseline (\S5.4 ``loop'').},label={lst:amber-loop}]
import ambr as am

class WealthAgent(am.Agent):
    def setup(self):
        self.wealth = 1

class WealthTransfer(am.Model):
    def setup(self):
        n = int(self.p["n"])
        self.add_agents(n, agent_class=WealthAgent)

    def step_oop(self):
        agents = list(self.agents)
        entry = {a.id: a.wealth for a in agents}
        delta = {a.id: 0 for a in agents}
        donors = [a for a in agents if entry[a.id] > 0]
        for donor in donors:
            recipient = self.random.choice(agents)
            delta[donor.id] -= 1
            delta[recipient.id] += 1
        for agent in agents:
            agent.wealth = entry[agent.id] + delta[agent.id]

WealthTransfer({"n": 10_000, "steps": 50}).cpu(mode="oop").run()
\end{lstlisting}

\paragraph{Vectorized CPU (\S5.4 ``vectorized'').}
The vectorized CPU lane computes donors and recipients from the entry column,
then applies one combined debit-and-credit \texttt{scatter\_add}. Contract
checking reports the operations it observes; it is not a
schedule-faithfulness proof.

\begin{lstlisting}[style=applisting,language=Python,caption={AMBER wealth transfer --- vectorized CPU idiom.},label={lst:amber-cpu}]
import ambr as am
import numpy as np

class WealthTransfer(am.Model):
    def setup(self):
        n = int(self.p["n"])
        self.add_agents(n, wealth=np.ones(n, dtype=np.int64))

    def step_vectorized(self):
        wealth = self.agents.wealth.to_numpy()
        donor_pos = np.flatnonzero(wealth > 0)
        if donor_pos.size == 0:
            return
        ids = self.agents.ids.to_numpy()
        recipients = self.rng.choice(ids, size=donor_pos.size)
        targets = np.concatenate((ids[donor_pos], recipients))
        deltas = np.concatenate((
            -np.ones(donor_pos.size, dtype=np.int64),
             np.ones(donor_pos.size, dtype=np.int64)))
        self.agents.at[targets].scatter_add(wealth=deltas)

WealthTransfer({"n": 10_000, "steps": 50}).cpu(
    mode="vectorized").run(contract="check")
\end{lstlisting}

\paragraph{GPU backend (\S5.4 ``GPU'').}
The benchmark calls
\texttt{model.approve\_fast\_path("workload smoke/invariant suite").gpu().run()}
to select a private no-contract loop. The approval string records a caller
decision; AMBER does not verify it. The fused kernel first computes
\texttt{donor\_positions = nonzero(wealth > 0)} and then atomically debits those
frozen donors and credits sampled recipients. The final ten-run mean is
94.3\,ms versus 193.5\,ms for FLAME GPU~2 (Table~\ref{tab:flame}). Because the
private loop bypasses the general instrumented runner, this timing is not
described as certified.

\appendixbreak
\subsection{Columnar Python: mesa-frames}
\label{app:usage:columnar}

mesa-frames \citep{mesaframes} stores agents in a Polars-backed
\texttt{AgentSetPolars}. Listing~\ref{lst:mesaframes} reads wealth to NumPy,
computes one snapshot result, and commits it with \texttt{set}. Section~5.1
shows that changing the local computation from one whole-array pass to
blockwise in-place mutation before the same final \texttt{set} changes the
operational schedule; this is an API-semantics observation, not a framework bug.
\begin{lstlisting}[style=applisting,language=Python,caption={mesa-frames wealth transfer (Polars \texttt{AgentSetPolars}).},label={lst:mesaframes}]
from mesa_frames import AgentSetPolars, ModelDF
import polars as pl
import numpy as np

class WealthSet(AgentSetPolars):
    def __init__(self, model, n):
        super().__init__(model)
        self.add(pl.DataFrame({
            "unique_id": np.arange(n, dtype=np.int64),
            "wealth": np.ones(n, dtype=np.int64),
        }))

    def step(self):
        ids = self.agents["unique_id"].to_numpy()
        wealth = self.agents["wealth"].to_numpy()
        donors = wealth > 0
        k = int(donors.sum())
        if k == 0:
            return
        recipients = self.model.rng.choice(ids, size=k)
        new = wealth.copy()
        pos = {int(u): i for i, u in enumerate(ids)}
        np.add.at(new, np.flatnonzero(donors), -1)
        np.add.at(new, [pos[int(r)] for r in recipients], 1)
        self.set("wealth", new)

class WealthModel(ModelDF):
    def __init__(self, n=10_000, steps=50):
        super().__init__(seed=0)
        self.rng = np.random.default_rng(0)
        self.steps = steps
        self.agents += WealthSet(self, n)

    def run(self):
        for _ in range(self.steps):
            self.agents.do("step")

WealthModel().run()
\end{lstlisting}

\appendixbreak
\subsection{Object-loop CPU frameworks}
\label{app:usage:loops}

These examples advance agents through scheduler-style APIs. The Python programs
incur per-agent interpreter work, while Agents.jl executes compiled Julia code.
In our released wealth implementations, the largest completed populations under
the fixed budget lie between $10^4$ and $10^6$ in the archived historical
scaling sweep; this is not asserted as an intrinsic framework limit.

\begin{lstlisting}[style=applisting,language=Python,caption={AgentPy wealth transfer.},label={lst:agentpy}]
import agentpy as ap

class WealthAgent(ap.Agent):
    def setup(self):
        self.wealth = 1

    def step(self):
        if self.wealth > 0:
            other = self.model.agents.random()
            if other is not self:
                self.wealth -= 1
                other.wealth += 1

class WealthModel(ap.Model):
    def setup(self):
        self.agents = ap.AgentList(self, self.p.n, WealthAgent)

    def step(self):
        self.agents.step()

WealthModel({"n": 10_000, "steps": 50}).run(steps=50, display=False)
\end{lstlisting}

\noindent\begin{minipage}{\textwidth}
\begin{lstlisting}[style=applisting,language=Python,caption={Mesa 3.x wealth transfer.},label={lst:mesa}]
import mesa

class WealthAgent(mesa.Agent):
    def __init__(self, model):
        super().__init__(model)
        self.wealth = 1

    def step(self):
        if self.wealth > 0:
            other = self.model.random.choice(list(self.model.agents))
            if other is not self:
                self.wealth -= 1
                other.wealth += 1

class WealthModel(mesa.Model):
    def __init__(self, n=10_000, seed=None):
        super().__init__(seed=seed)
        for _ in range(n):
            WealthAgent(self)

    def step(self):
        self.agents.shuffle_do("step")

m = WealthModel()
for _ in range(50):
    m.step()
\end{lstlisting}
\end{minipage}

\begin{lstlisting}[style=applisting,language=Julia,caption={Agents.jl wealth transfer.},label={lst:agentsjl}]
using Agents

@agent struct WealthAgent(NoSpaceAgent)
    wealth::Float64
end

function wealth_step!(agent, model)
    agent.wealth > 0 || return
    partner = random_agent(model)
    if partner !== nothing && partner.id != agent.id
        agent.wealth -= 1
        partner.wealth += 1
    end
end

model = StandardABM(WealthAgent; agent_step! = wealth_step!)
for _ in 1:10_000
    add_agent!(model, 1.0)
end
step!(model, 50)
\end{lstlisting}

\appendixbreak
\subsection{FLAME GPU~2}
\label{app:usage:flame}

FLAME GPU~2 \citep{richmond2023} is the GPU baseline in \S5.4. It separates
\emph{host} model construction (Python) from \emph{device} agent logic
(runtime-compiled CUDA-C). Wealth transfer uses two RTC layers:
\texttt{wt\_give} debits donors and posts a bucket message keyed by recipient
ID; \texttt{wt\_receive} credits the recipient. Listing~\ref{lst:flame} shows
the give function; the host wires the layers, allocates an
\texttt{AgentVector}, and calls the simulation (full host code is in the
archive). The corrected ten-run campaign finds AMBER faster in all four rows of
Table~\ref{tab:flame}, but Schelling includes Python per-agent FLAME
initialization in the timed region and is exploratory. The transition details
and random-number machinery are not identical, so these results compare our
implementations rather than establishing universal framework rankings. The historical SLOC artifact also
excludes AMBER's private backend and is not used as a conciseness claim.
\begin{lstlisting}[style=applisting,language=Python,caption={FLAME GPU~2 --- RTC give layer (wealth transfer).},label={lst:flame}]
import pyflamegpu as fg

WT_GIVE = r"""
FLAMEGPU_AGENT_FUNCTION(wt_give, flamegpu::MessageNone,
                        flamegpu::MessageBucket) {
    const int w = FLAMEGPU->getVariable<int>("wealth");
    if (w > 0) {
        const unsigned int n =
            FLAMEGPU->environment.getProperty<unsigned int>("n");
        const unsigned int r =
            FLAMEGPU->random.uniform<unsigned int>(1u, n);
        FLAMEGPU->message_out.setKey(r);
        FLAMEGPU->message_out.setVariable<int>("amount", 1);
        FLAMEGPU->setVariable<int>("wealth", w - 1);
    }
    return flamegpu::ALIVE;
}"""
\end{lstlisting}

\end{document}
